\subsection{Supporting Analysis of Executed Behavior}\label{sec:supporting}\label{sec:f1-collapse}
The MATH controls show that outcome differences need not require
different programs. The supporting BIRD study asks the converse:
do source differences produce distinct, useful execution? Here a
mechanism is an implemented prompt transformation, tool use, or
control-flow operation. This study uses separate populations, judging,
and execution protocols; it does not explain the MATH results causally.

\paragraph{Source differences and observed outcomes.}
All $12$ discovery candidates pass a code-difference check, yet $21/66$
generated-candidate pairs have identical correctness vectors
($36/91$ including baselines). Source similarity has essentially no
rank association with correctness disagreement ($\rho=-0.010$).
Agreement on these $60$ tasks is not general program equivalence,
but motivates inspecting execution (Appendix Figure~\ref{fig:phenomenon}).

\paragraph{Traces locate three distinct failures.}
\textbf{T1: described but absent.} Six of twelve candidates claim
retrieval or self-checking but make one model call and execute no
artifact. \textbf{T2: present but untriggered.} On all ten audited
\react{} tasks, the first SQL executes successfully, so the repair
branch never fires and final SQL equals \bare{}; execution success
is not semantic correctness. \textbf{T3: executed but broken.} One
candidate runs multi-turn repair but returns prose as final SQL,
losing tasks that the baseline answers correctly. These failures
require implementation checks, branch-exercising tasks, and valid
output contracts, respectively. They do not establish that activation
improves accuracy. Appendix~\ref{app:bird-traces} retains worked traces,
hand-written controls, and their development-exposure limitations.

The separate BIRD admission study also finds a mean headroom contrast
of $\RevisionPrimaryDelta$\,\% for forced/gated versus free/ungated
generation. Its gate bundles structural exclusion with conformance
verification, so it does not identify a pure verification effect
(Appendix~\ref{app:phase2-protocol}). Source and admission evidence
therefore require an execution check; persistent task advantages and
selection value remain separate measurements.
